\documentclass[a4paper,12pt]{article}

\usepackage[utf8]{inputenc}
\usepackage[T1]{fontenc}
\usepackage[italian]{babel}
\usepackage{geometry}
\usepackage{hyperref}
\usepackage{listings}
\usepackage{xcolor}
\usepackage{amsmath}
\usepackage{enumitem}

\geometry{margin=2.5cm}

\definecolor{codegray}{rgb}{0.95,0.95,0.95}
\lstset{
  basicstyle=\ttfamily\small,
  backgroundcolor=\color{codegray},
  frame=single,
  breaklines=true,
  columns=fullflexible,
  xleftmargin=\parindent,
  xrightmargin=\parindent
}

\title{Relazione Tecnica sullo Sviluppo di una Binary Heap Generica in C}
\author{}
\date{\today}

\begin{document}

\maketitle

\tableofcontents
\newpage

\section{Introduzione}

L'obiettivo del lavoro svolto in questa directory è stato progettare e implementare una struttura dati di tipo \emph{binary heap} generica in linguaggio C, prendendo come punto di partenza un'analisi dei sorgenti della \texttt{glibc} 2.36 presenti nel progetto tramite symlink. Il risultato finale comprende:

\begin{itemize}[leftmargin=*]
    \item un'implementazione generica della heap basata su \texttt{void **} e funzione di confronto;
    \item un file header con l'interfaccia pubblica;
    \item un programma dimostrativo eseguibile;
    \item un file di test separato con controlli automatici;
    \item un \texttt{Makefile} per build, esecuzione della demo e lancio dei test.
\end{itemize}

La richiesta iniziale non era soltanto ``scrivere una heap'', ma capire prima se all'interno della \texttt{glibc} fosse già presente una struttura dati simile o riusabile. Per questo motivo il lavoro è partito da una fase di esplorazione del codice sorgente, seguita poi da una fase di progettazione e infine dalla realizzazione vera e propria.

\section{Analisi Preliminare dei Sorgenti glibc}

\subsection{Accesso ai sorgenti}

All'interno della directory di lavoro era presente un symlink chiamato \texttt{glibc-2.36}, puntante ai sorgenti reali della libreria GNU C. La prima attività è consistita nel verificare:

\begin{itemize}[leftmargin=*]
    \item che il symlink fosse leggibile;
    \item che il target fosse effettivamente accessibile dal workspace;
    \item che la directory contenesse i file tipici dei sorgenti \texttt{glibc}, come \texttt{configure}, \texttt{Makerules}, \texttt{README} e \texttt{COPYING}.
\end{itemize}

La verifica ha confermato che il sorgente era disponibile e ispezionabile.

\subsection{Ricerca di una binary heap o priority queue già presente}

La seconda fase è stata una ricerca mirata nei sorgenti \texttt{glibc} per stabilire se esistesse già:

\begin{itemize}[leftmargin=*]
    \item una struttura dati dichiarata esplicitamente come \emph{binary heap};
    \item una \emph{priority queue} riusabile;
    \item una utility generica che implementasse operazioni equivalenti;
    \item un'implementazione locale che potesse essere estratta o adattata.
\end{itemize}

La ricerca è stata condotta usando parole chiave come:

\begin{itemize}[leftmargin=*]
    \item \texttt{binary heap};
    \item \texttt{priority queue}, \texttt{priority\_queue}, \texttt{pqueue}, \texttt{pq};
    \item \texttt{heapify}, \texttt{sift\_down}, \texttt{sift\_up};
    \item nomi file contenenti \texttt{heap}, \texttt{priority}, \texttt{queue}.
\end{itemize}

\subsection{Esito della ricerca}

La conclusione dell'analisi è stata netta:

\begin{itemize}[leftmargin=*]
    \item nei sorgenti \texttt{glibc-2.36} non è presente una \emph{priority queue} generica pronta all'uso;
    \item non è presente una struttura dati esposta o documentata come \emph{binary heap} riusabile;
    \item esistono però due elementi interessanti:
          \begin{itemize}
              \item una coda ordinata per timer, implementata come lista ordinata;
              \item un'implementazione locale di \emph{heapsort} che utilizza esplicitamente la proprietà di heap binario su array.
          \end{itemize}
\end{itemize}

\subsection{La coda ordinata dei timer}

Nel file \texttt{fbtl/sysdeps/pthread/timer\_routines.c} è presente una struttura che, a prima vista, ricorda una priority queue: i timer vengono inseriti in ordine di scadenza. Tuttavia l'implementazione non usa un heap, ma una lista doppiamente collegata ordinata. L'inserimento avviene tramite scansione lineare fino al punto corretto.

Questo approccio ha caratteristiche diverse rispetto a una binary heap:

\begin{itemize}[leftmargin=*]
    \item inserimento in tempo lineare \(O(n)\);
    \item accesso al minimo in tempo costante \(O(1)\), essendo in testa;
    \item nessuna rappresentazione su array;
    \item nessun \emph{sift up} o \emph{sift down}.
\end{itemize}

Per questo motivo tale codice non era un punto di partenza adeguato per implementare una heap binaria vera e propria.

\subsection{Il riferimento utile: \texttt{frame\_heapsort}}

Nel file \texttt{sysdeps/generic/unwind-dw2-fde.c} è stata individuata la funzione \texttt{frame\_heapsort}, che usa chiaramente la rappresentazione classica di un heap binario su array:

\begin{itemize}[leftmargin=*]
    \item figlio sinistro all'indice \(2i + 1\);
    \item figlio destro all'indice \(2i + 2\);
    \item ripristino della proprietà di heap tramite scambi successivi;
    \item selezione del figlio con priorità maggiore per mantenere l'invariante.
\end{itemize}

Pur essendo specializzata per il contesto interno della \texttt{glibc} e per una funzione di ordinamento, questa implementazione conteneva l'idea algoritmica giusta per derivare una struttura dati generica.

\section{Scelte Progettuali}

\subsection{Obiettivo della nuova implementazione}

L'obiettivo è stato definire una heap riusabile in C con queste proprietà:

\begin{itemize}[leftmargin=*]
    \item genericità sugli elementi memorizzati;
    \item separazione tra struttura dati e criterio di ordinamento;
    \item API semplice e minimale;
    \item dipendenze standard minime;
    \item integrazione facile con il \texttt{Makefile}.
\end{itemize}

\subsection{Rappresentazione dei dati}

La heap è stata implementata come array dinamico di puntatori generici:

\begin{lstlisting}[language=C]
struct binaryheap
{
    void **data;
    size_t size;
    size_t capacity;
    binaryheap_cmp_fn cmp;
};
\end{lstlisting}

Questa scelta comporta diversi vantaggi:

\begin{itemize}[leftmargin=*]
    \item la struttura non copia gli oggetti, ma memorizza puntatori;
    \item può contenere elementi di qualsiasi tipo;
    \item il criterio di priorità è delegato a una funzione di confronto;
    \item l'allocazione e la crescita della memoria sono gestite in modo centralizzato.
\end{itemize}

\subsection{Convenzione del comparatore}

Il comparatore segue una convenzione in stile \texttt{qsort}:

\begin{itemize}[leftmargin=*]
    \item ritorno \(< 0\): il primo argomento ha priorità più alta;
    \item ritorno \(= 0\): i due elementi sono equivalenti;
    \item ritorno \(> 0\): il secondo argomento ha priorità più alta.
\end{itemize}

Con questa convenzione, l'implementazione realizzata si comporta come una \emph{min-heap}: l'elemento ``più piccolò' secondo il comparatore risale verso la radice.

\subsection{Perchè una min-heap}

La scelta di implementare una min-heap e non una max-heap deriva da motivi pragmatici:

\begin{itemize}[leftmargin=*]
    \item molti usi tipici della priority queue richiedono l'estrazione del minimo;
    \item il comportamento è naturale se il comparatore ordina in senso crescente;
    \item passare da min-heap a max-heap è comunque banale: basta invertire la logica del comparatore.
\end{itemize}

\section{Implementazione del File \texttt{binaryheap.c}}

\subsection{Interfaccia pubblica}

Per rendere l'implementazione usabile sia dal programma demo sia dal file di test, è stato creato anche il file \texttt{binaryheap.h}, che espone:

\begin{itemize}[leftmargin=*]
    \item il tipo funzione \texttt{binaryheap\_cmp\_fn};
    \item la struttura \texttt{struct binaryheap};
    \item le funzioni pubbliche della heap.
\end{itemize}

Le operazioni esposte sono:

\begin{itemize}[leftmargin=*]
    \item \texttt{binaryheap\_init};
    \item \texttt{binaryheap\_destroy};
    \item \texttt{binaryheap\_reserve};
    \item \texttt{binaryheap\_push};
    \item \texttt{binaryheap\_peek};
    \item \texttt{binaryheap\_pop};
    \item \texttt{binaryheap\_size};
    \item \texttt{binaryheap\_empty};
    \item \texttt{binaryheap\_build}.
\end{itemize}

\subsection{\texttt{binaryheap\_init}}

Questa funzione inizializza la struttura in uno stato vuoto:

\begin{itemize}[leftmargin=*]
    \item puntatore dati a \texttt{NULL};
    \item dimensione a zero;
    \item capacità a zero;
    \item salvataggio del comparatore.
\end{itemize}

L'inizializzazione non alloca memoria subito. La memoria viene richiesta solo al primo inserimento o in seguito a una \texttt{reserve}.

\subsection{\texttt{binaryheap\_destroy}}

Questa funzione libera la memoria allocata dinamicamente per l'array interno. Non dealloca gli elementi puntati, perchè la heap non ne è proprietaria. Questa è una scelta intenzionale: la struttura gestisce i puntatori, non il ciclo di vita degli oggetti referenziati.

\subsection{\texttt{binaryheap\_reserve}}

Questa funzione permette di preallocare memoria per almeno un certo numero di elementi. Se la capacità richiesta è inferiore o uguale a quella già disponibile, non viene eseguita alcuna riallocazione.

Se invece serve aumentare la memoria:

\begin{itemize}[leftmargin=*]
    \item viene invocata \texttt{realloc};
    \item in caso di successo si aggiorna la capacità;
    \item in caso di fallimento si ritorna \texttt{-1}.
\end{itemize}

La funzione è utile sia per ottimizzazione, sia come building block per \texttt{push}.

\subsection{\texttt{binaryheap\_push}}

L'inserimento di un elemento segue due passi:

\begin{enumerate}[leftmargin=*]
    \item se necessario, espande la capacità dell'array;
    \item inserisce il nuovo elemento in fondo e lo fa risalire tramite \emph{sift up}.
\end{enumerate}

La crescita della capacità segue una politica geometrica:

\begin{itemize}[leftmargin=*]
    \item da \(0\) a \(8\) elementi alla prima espansione;
    \item successivamente raddoppio della capacità.
\end{itemize}

Questa scelta riduce il numero di riallocazioni e fornisce costo ammortizzato efficiente sugli inserimenti.

\subsection{\texttt{binaryheap\_peek}}

Restituisce l'elemento in radice senza rimuoverlo. Se la heap è vuota, ritorna \texttt{NULL}. In una min-heap questo elemento corrisponde al minimo secondo il comparatore.

\subsection{\texttt{binaryheap\_pop}}

L'estrazione dell'elemento in testa avviene come segue:

\begin{enumerate}[leftmargin=*]
    \item salva l'elemento in radice;
    \item sposta l'ultimo elemento della heap in posizione \(0\);
    \item decrementa la dimensione;
    \item ripristina la proprietà di heap con \emph{sift down}.
\end{enumerate}

Se la heap è vuota, ritorna \texttt{NULL}.

\subsection{\texttt{binaryheap\_size} e \texttt{binaryheap\_empty}}

Queste funzioni forniscono introspezione minima sulla struttura:

\begin{itemize}[leftmargin=*]
    \item \texttt{binaryheap\_size} restituisce il numero di elementi;
    \item \texttt{binaryheap\_empty} ritorna vero se la heap è vuota.
\end{itemize}

\subsection{\texttt{binaryheap\_build}}

Questa funzione costruisce una heap da un array di puntatori già esistente. Invece di inserire gli elementi uno a uno, applica una costruzione bottom-up:

\begin{itemize}[leftmargin=*]
    \item inizializza i campi della struttura;
    \item parte dall'ultimo nodo interno;
    \item applica \emph{sift down} fino alla radice.
\end{itemize}

Questa tecnica permette una costruzione in tempo lineare \(O(n)\), più efficiente rispetto a \(n\) inserimenti consecutivi.

\section{Dettagli Algoritmici}

\subsection{\texttt{binaryheap\_sift\_up}}

L'operazione di risalita confronta un nodo con il padre:

\begin{itemize}[leftmargin=*]
    \item se il nodo ha priorità più alta del padre, avviene uno scambio;
    \item il processo continua finchè l'invariante di heap non è ristabilito o si raggiunge la radice.
\end{itemize}

Nel caso peggiore, il numero di passi è proporzionale all'altezza dell'albero, quindi \(O(\log n)\).

\subsection{\texttt{binaryheap\_sift\_down}}

L'operazione di discesa confronta un nodo con i figli:

\begin{itemize}[leftmargin=*]
    \item viene scelto il figlio con priorità maggiore tra i due figli disponibili;
    \item se tale figlio ha priorità maggiore del nodo corrente, i due vengono scambiati;
    \item l'algoritmo continua dal nuovo indice.
\end{itemize}

Anche in questo caso il costo nel caso peggiore è \(O(\log n)\).

\subsection{Invariante della struttura}

L'invariante fondamentale della heap è:

\[
    \text{per ogni indice } i, \quad data[i] \leq data[2i+1] \text{ e } data[i] \leq data[2i+2]
\]

naturalmente nel senso definito dal comparatore.

Questo garantisce che la radice contenga sempre l'elemento con priorità più alta nel senso scelto dalla struttura, cioè il minimo nel caso corrente.

\section{Complessità Computazionale}

Le complessità delle operazioni principali sono:

\begin{center}
    \begin{tabular}{|l|c|}
        \hline
        Operazione                   & Complessità                \\
        \hline
        \texttt{binaryheap\_init}    & \(O(1)\)                   \\
        \texttt{binaryheap\_destroy} & \(O(1)\)                   \\
        \texttt{binaryheap\_peek}    & \(O(1)\)                   \\
        \texttt{binaryheap\_empty}   & \(O(1)\)                   \\
        \texttt{binaryheap\_size}    & \(O(1)\)                   \\
        \texttt{binaryheap\_push}    & \(O(\log n)\) ammortizzato \\
        \texttt{binaryheap\_pop}     & \(O(\log n)\)              \\
        \texttt{binaryheap\_build}   & \(O(n)\)                   \\
        \hline
    \end{tabular}
\end{center}

L'uso di un array dinamico rende inoltre l'accesso ai nodi molto efficiente dal punto di vista cache, rispetto a strutture basate su puntatori espliciti a figli e genitori.

\section{Programma Demo}

\subsection{Scopo di \texttt{binaryheap\_demo.c}}

È stato creato un programma dimostrativo, \texttt{binaryheap\_demo.c}, per mostrare l'uso base della struttura:

\begin{itemize}[leftmargin=*]
    \item definizione di un comparatore per interi;
    \item inizializzazione della heap;
    \item inserimento di alcuni valori;
    \item estrazione in ordine crescente;
    \item stampa dei risultati.
\end{itemize}

Questo file non nasce come test formale, ma come esempio d'uso e verifica manuale rapida.

\subsection{Differenza tra demo e test}

La demo e il test hanno scopi diversi:

\begin{itemize}[leftmargin=*]
    \item la demo mostra \emph{come usare} l'API;
    \item il test controlla \emph{se l'API funziona correttamente}.
\end{itemize}

Per questo motivo la demo è stata lasciata separata dal file di test.

\section{Testing}

\subsection{Motivazione}

Dopo la demo, è emersa la necessità di avere un file di test separato. Il motivo principale è che una demo può dare l'impressione che il codice funzioni, ma non fornisce garanzie sistematiche. Un test dedicato invece:

\begin{itemize}[leftmargin=*]
    \item verifica più scenari;
    \item fallisce automaticamente in presenza di errori;
    \item è integrabile in un workflow di build;
    \item può essere esteso facilmente.
\end{itemize}

\subsection{File \texttt{test\_binaryheap.c}}

È stato quindi aggiunto il file \texttt{test\_binaryheap.c}, costruito con \texttt{assert}. I casi coperti sono:

\begin{itemize}[leftmargin=*]
    \item \textbf{heap vuota}: verifica di \texttt{empty}, \texttt{size}, \texttt{peek} e \texttt{pop} su struttura inizializzata ma senza elementi;
    \item \textbf{ordine di inserimento/estrazione}: inserimento di una sequenza non ordinata e controllo dell'estrazione in ordine crescente;
    \item \textbf{duplicati}: verifica del corretto comportamento in presenza di valori ripetuti;
    \item \textbf{build}: costruzione di una heap a partire da un array già esistente e verifica dell'ordine finale.
\end{itemize}

\subsection{Output dei test}

Inizialmente il file di test non produceva output in caso di successo, comportamento corretto ma poco evidente a schermo. Successivamente è stato aggiunto un messaggio finale:

\begin{lstlisting}
All tests passed
\end{lstlisting}

In questo modo:

\begin{itemize}[leftmargin=*]
    \item se un \texttt{assert} fallisce, il programma termina con errore;
    \item se tutto va bene, il test stampa esplicitamente il successo.
\end{itemize}

\section{Build Automation con Makefile}

\subsection{Obiettivi del Makefile}

Il \texttt{Makefile} è stato creato per gestire in modo semplice e ripetibile le attività essenziali:

\begin{itemize}[leftmargin=*]
    \item compilazione del file oggetto;
    \item compilazione ed esecuzione della demo;
    \item compilazione ed esecuzione dei test;
    \item pulizia dei file generati.
\end{itemize}

\subsection{Target disponibili}

I target principali sono:

\begin{itemize}[leftmargin=*]
    \item \texttt{make}: compila \texttt{binaryheap.o};
    \item \texttt{make run}: compila ed esegue \texttt{binaryheap\_demo};
    \item \texttt{make test}: compila ed esegue \texttt{test\_binaryheap};
    \item \texttt{make clean}: rimuove artefatti di compilazione ed eseguibili.
\end{itemize}

\subsection{Scelta delle flag di compilazione}

Le flag configurate sono:

\begin{lstlisting}
-std=c11 -Wall -Wextra -pedantic
\end{lstlisting}

Su richiesta esplicita, non sono state abilitate ottimizzazioni del compilatore. Questa scelta è utile in una fase didattica o di sviluppo iniziale per diversi motivi:

\begin{itemize}[leftmargin=*]
    \item mantiene il comportamento più prevedibile;
    \item facilita il debugging;
    \item evita trasformazioni aggressive che potrebbero complicare l'analisi del codice generato.
\end{itemize}

\section{Struttura Finale dei File}

Al termine del lavoro, i file principali del progetto risultano essere:

\begin{itemize}[leftmargin=*]
    \item \texttt{binaryheap.c}: implementazione della struttura dati;
    \item \texttt{binaryheap.h}: interfaccia pubblica;
    \item \texttt{binaryheap\_demo.c}: programma dimostrativo;
    \item \texttt{test\_binaryheap.c}: file di test automatici;
    \item \texttt{Makefile}: build automation;
    \item \texttt{report/report.tex}: questa relazione.
\end{itemize}

\section{Valutazione del Lavoro Svolto}

\subsection{Risultati raggiunti}

Il lavoro ha prodotto una struttura dati funzionante, generica e testata, costruita a partire da un'analisi concreta dei sorgenti \texttt{glibc}. Gli obiettivi principali sono stati raggiunti:

\begin{itemize}[leftmargin=*]
    \item identificazione dell'assenza di una priority queue generica in \texttt{glibc};
    \item individuazione di un riferimento algoritmico utile;
    \item implementazione di una min-heap generica in C;
    \item predisposizione di un ambiente minimo di build e test.
\end{itemize}

\subsection{Limiti attuali}

La soluzione attuale è adeguata per lo scopo richiesto, ma presenta alcuni limiti naturali:

\begin{itemize}[leftmargin=*]
    \item non gestisce direttamente la deallocazione degli oggetti memorizzati;
    \item non include operazioni come \texttt{decrease-key}, \texttt{remove} arbitrario o aggiornamento di priorità;
    \item non è presente una suite di test ancora più estesa con scenari stressati o casuali;
    \item non è stata prodotta una libreria statica o condivisa separata.
\end{itemize}

\subsection{Possibili estensioni}

Tra le evoluzioni più naturali del progetto si possono indicare:

\begin{itemize}[leftmargin=*]
    \item aggiunta di test con sequenze miste di \texttt{push} e \texttt{pop};
    \item supporto a max-heap tramite wrapper o comparatori dedicati;
    \item introduzione di un indice esterno per supportare \texttt{decrease-key};
    \item produzione di una libreria \texttt{.a};
    \item documentazione aggiuntiva dell'API in formato header commentato o manuale.
\end{itemize}

\section{Conclusione}

L'attività svolta ha seguito un percorso tecnico coerente:

\begin{enumerate}[leftmargin=*]
    \item verifica della disponibilità dei sorgenti \texttt{glibc};
    \item analisi per cercare strutture dati equivalenti a una priority queue o binary heap;
    \item identificazione di \texttt{frame\_heapsort} come miglior riferimento algoritmico;
    \item progettazione di una heap generica basata su \texttt{void **} e comparatore;
    \item realizzazione dell'implementazione e della relativa interfaccia;
    \item aggiunta di demo, test e automazione della build;
    \item stesura della presente relazione.
\end{enumerate}

Il risultato è un piccolo progetto C pulito e coerente, nato da un'esigenza pratica e sviluppato con un approccio incrementale: prima l'analisi, poi la progettazione, infine implementazione e verifica.

\end{document}
